% ===========================================================================
%  preamble.tex — Cấu hình gói, font, kiểu trình bày cho báo cáo Synergit
%  Biên dịch bằng XeLaTeX / Tectonic (hỗ trợ tiếng Việt qua fontspec).
% ===========================================================================

\usepackage{fontspec}
% Dùng font mặc định của LaTeX: Latin Modern (Computer Modern) — hỗ trợ đầy đủ
% glyph tiếng Việt. fontspec tự nạp Latin Modern khi không chỉ định font khác.
\setmainfont{lmroman10-regular.otf}[
  BoldFont = lmroman10-bold.otf,
  ItalicFont = lmroman10-italic.otf,
  BoldItalicFont = lmroman10-bolditalic.otf ]
\setsansfont{lmsans10-regular.otf}[
  BoldFont = lmsans10-bold.otf,
  ItalicFont = lmsans10-oblique.otf,
  BoldItalicFont = lmsans10-boldoblique.otf ]
\setmonofont{lmmonolt10-regular.otf}[
  BoldFont = lmmonolt10-bold.otf,
  ItalicFont = lmmonolt10-oblique.otf,
  BoldItalicFont = lmmonolt10-boldoblique.otf, Scale = 0.95 ]

\usepackage{geometry}
\geometry{a4paper, left=3cm, right=2cm, top=2cm, bottom=2cm}

\usepackage{amsmath,amssymb}
\usepackage{graphicx}
\usepackage{xcolor}
\usepackage{colortbl}
\usepackage{array}
\usepackage{longtable}
\usepackage{booktabs}
\usepackage{tabularx}
\usepackage{multirow}
\usepackage{enumitem}
\usepackage{ragged2e}
\usepackage{float}
\usepackage{caption}
\usepackage{titlesec}
\usepackage{titletoc}
\usepackage{fancyhdr}
\usepackage{setspace}
\usepackage{parskip}
\usepackage[hidelinks]{hyperref}
\usepackage{xurl}   % cho phép URL dài tự ngắt dòng
\usepackage{tikz}
\usetikzlibrary{shapes.geometric, arrows.meta, positioning, fit, calc, backgrounds, matrix}
\usepackage{tcolorbox}

% --- Màu chủ đề (Synergit: xanh lam GitHub-ish) -----------------------------
\definecolor{sgBlue}{HTML}{0969DA}
\definecolor{sgInk}{HTML}{1F2328}
\definecolor{sgCream}{HTML}{F6F8FA}
\definecolor{sgLine}{HTML}{D0D7DE}
\definecolor{sgAccent}{HTML}{2DA44E}
\colorlet{bkCream}{sgCream}
\colorlet{bkLine}{sgLine}
\colorlet{bkBrown}{sgInk}
\definecolor{codebg}{HTML}{F4F4F4}
\definecolor{codekw}{HTML}{0969DA}
\definecolor{codecm}{HTML}{6A737D}
\definecolor{codestr}{HTML}{1B7A43}

% Kiểu dùng chung cho các sơ đồ ERD theo schema
\tikzset{
  erdent/.style={draw=sgBlue, fill=sgCream, rounded corners, align=left,
                 inner sep=4pt, font=\scriptsize, text width=3.05cm},
  erdext/.style={draw=sgLine, dashed, fill=white, rounded corners, align=center,
                 inner sep=4pt, font=\scriptsize\itshape, text width=2.6cm},
  erdrel/.style={draw=sgLine, thick},
  erdsoft/.style={draw=sgLine, thick, dashed},
  erdcard/.style={font=\scriptsize, text=sgBlue, fill=white, inner sep=1pt},
}
\tcbuselibrary{skins, breakable}
\usepackage{listings}

% --- Cỡ chữ 13pt, giãn dòng 1.5 (chuẩn báo cáo đồ án) -----------------------
\renewcommand{\normalsize}{\fontsize{13pt}{19.5pt}\selectfont}
\normalsize
\setstretch{1.42}

% --- Canh lề hai bên, hạn chế ngắt từ tiếng Việt ---------------------------
\justifying
\hyphenpenalty=8000
\tolerance=2000
\setlength{\emergencystretch}{3em}

% --- Nhãn tiếng Việt cho Hình / Bảng / Mục lục -----------------------------
\captionsetup{labelsep=colon, font=small, labelfont=bf}
\captionsetup[figure]{name=Hình}
\captionsetup[table]{name=Bảng}
\renewcommand{\chaptername}{Chương}
\renewcommand{\contentsname}{MỤC LỤC}
\renewcommand{\listfigurename}{MỤC LỤC HÌNH ẢNH}
\renewcommand{\listtablename}{MỤC LỤC BẢNG}
\renewcommand{\bibname}{TÀI LIỆU THAM KHẢO}
\renewcommand{\figurename}{Hình}
\renewcommand{\tablename}{Bảng}

% --- Đánh số Hình/Bảng theo chương -----------------------------------------
\usepackage{chngcntr}
\counterwithin{figure}{chapter}
\counterwithin{table}{chapter}

% --- Kiểu tiêu đề Chương / Mục ---------------------------------------------
\titleformat{\chapter}[display]
  {\centering\bfseries\fontsize{15pt}{18pt}\selectfont}
  {\MakeUppercase{\chaptertitlename}\ \thechapter}
  {6pt}
  {\MakeUppercase}
\titlespacing*{\chapter}{0pt}{0pt}{18pt}

\titleformat{\section}
  {\bfseries\fontsize{13.5pt}{16pt}\selectfont}{\thesection.}{0.5em}{}
\titleformat{\subsection}
  {\bfseries\itshape\fontsize{13pt}{16pt}\selectfont}{\thesubsection.}{0.5em}{}
\titleformat{\subsubsection}
  {\itshape\fontsize{13pt}{16pt}\selectfont}{\thesubsubsection.}{0.5em}{}
\titlespacing*{\section}{0pt}{10pt}{4pt}
\titlespacing*{\subsection}{0pt}{8pt}{3pt}
\titlespacing*{\subsubsection}{0pt}{6pt}{3pt}

% --- Header / footer (số trang góc dưới) -----------------------------------
\pagestyle{fancy}
\fancyhf{}
\renewcommand{\headrulewidth}{0pt}
\fancyfoot[C]{\thepage}
\fancypagestyle{plain}{\fancyhf{}\fancyfoot[C]{\thepage}\renewcommand{\headrulewidth}{0pt}}

% --- Danh sách gọn hơn ------------------------------------------------------
\setlist[itemize]{leftmargin=1.4em, topsep=2pt, itemsep=1pt, parsep=0pt}
\setlist[enumerate]{leftmargin=1.8em, topsep=2pt, itemsep=1pt, parsep=0pt}

% --- Kiểu khối mã nguồn -----------------------------------------------------
\lstdefinestyle{sg}{
  backgroundcolor=\color{codebg},
  basicstyle=\ttfamily\footnotesize,
  keywordstyle=\color{codekw}\bfseries,
  commentstyle=\color{codecm}\itshape,
  stringstyle=\color{codestr},
  numbers=left, numberstyle=\tiny\color{codecm}, numbersep=8pt,
  breaklines=true, showstringspaces=false, tabsize=2,
  frame=single, rulecolor=\color{sgLine},
  framexleftmargin=4pt, xleftmargin=14pt, aboveskip=8pt, belowskip=6pt,
  columns=fullflexible, keepspaces=true,
}
\lstset{style=sg}
\lstdefinelanguage{Go}{
  keywords={func,package,import,type,struct,interface,map,chan,go,defer,return,if,else,for,range,switch,case,default,var,const,nil,error,select},
  ndkeywords={string,int,int64,bool,uint,byte,UUID,context},
  sensitive=true, comment=[l]{//}, morecomment=[s]{/*}{*/},
  morestring=[b]", morestring=[b]`,
}
\lstdefinelanguage{TypeScript}{
  keywords={const,let,var,function,async,await,return,if,else,for,of,in,while,
            import,export,from,type,interface,class,extends,implements,new,
            public,private,readonly,as,enum,void,null,undefined,true,false},
  ndkeywords={string,number,boolean,Promise,Record,Array,Partial},
  sensitive=true, comment=[l]{//}, morecomment=[s]{/*}{*/},
  morestring=[b]", morestring=[b]', morestring=[b]`,
}

% --- Lệnh chèn ảnh placeholder (chờ ảnh chụp màn hình thật) -----------------
% \phimg{<caption>}{<chiều cao>}  — ví dụ \phimg{Giao diện trang Home}{6cm}
\newcommand{\phimg}[2]{%
  \begin{figure}[H]\centering
  \begin{tcolorbox}[width=0.88\linewidth, height=#2, halign=center, valign=center,
      colback=sgCream, colframe=sgLine, sharp corners, boxrule=0.6pt,
      before skip=4pt, after skip=2pt]
      {\color{sgInk}\itshape [\,Ảnh minh hoạ:\ #1\,]\\[4pt]\footnotesize (chèn ảnh chụp màn hình thật tại đây)}
  \end{tcolorbox}
  \caption{#1}
  \end{figure}}

% --- Lệnh chèn ảnh "thông minh" (skeleton) ---------------------------------
% \figauto[<nhãn>]{<đường dẫn file>}{<caption>}{<chiều cao tối đa>}
%   • Nếu file ảnh TỒN TẠI  -> tự chèn ảnh (giới hạn theo bề rộng & chiều cao).
%   • Nếu CHƯA có file       -> hiện khung placeholder ghi rõ tên file cần bỏ vào.
\newcommand{\figautolabel}{}
\newcommand{\figauto}[4][]{%
  \renewcommand{\figautolabel}{#1}%
  \begin{figure}[H]\centering
  \IfFileExists{#2}{%
    \includegraphics[width=0.95\linewidth, height=\dimexpr#4*13/10\relax, keepaspectratio]{#2}%
  }{%
    \begin{tcolorbox}[width=0.88\linewidth, height=#4, halign=center, valign=center,
        colback=sgCream, colframe=sgLine, sharp corners, boxrule=0.6pt,
        before skip=4pt, after skip=2pt]
        {\color{sgInk}\itshape [\,Đặt ảnh vào:\ \texttt{#2}\,]\\[4pt]
         \footnotesize #3}
    \end{tcolorbox}%
  }%
  \caption{#3}%
  \ifx\figautolabel\empty\else\label{#1}\fi
  \end{figure}}

% --- Tiêu đề trang phụ (không đánh số chương) cho phần đầu ------------------
\newcommand{\plainheading}[1]{%
  \par\vspace{6pt}\begin{center}\bfseries\fontsize{15pt}{18pt}\selectfont\MakeUppercase{#1}\end{center}\vspace{4pt}}

% --- Bảng đặc tả Usecase ----------------------------------------------------
% Dùng:  \begin{ucspec}{nhãn}{Tiêu đề bảng}  ... các dòng "Trường & Giá trị \\\hline" ... \end{ucspec}
\newenvironment{ucspec}[2]{%
  \renewcommand{\arraystretch}{1.25}\small
  \begin{longtable}{|>{\bfseries\raggedright\arraybackslash}p{3.4cm}|p{10.3cm}|}
  \caption{#2}\label{#1}\\
  \hline\endfirsthead
  \hline\multicolumn{2}{|r|}{\textit{(tiếp theo)}}\\\hline\endhead
}{%
  \end{longtable}\normalsize}

% --- Bảng đặc tả chi tiết bảng dữ liệu --------------------------------------
% Dùng:  \begin{dbtable}{nhãn}{Tiêu đề}  ... "STT & trường & kiểu & mô tả \\\hline" ... \end{dbtable}
\newenvironment{dbtable}[2]{%
  \renewcommand{\arraystretch}{1.2}\small
  \begin{longtable}{|c|>{\ttfamily}p{3.5cm}|>{\ttfamily}p{2.8cm}|p{5.3cm}|}
  \caption{#2}\label{#1}\\
  \hline\rowcolor{sgCream}
  \textbf{STT} & \normalfont\textbf{Tên trường} & \normalfont\textbf{Kiểu dữ liệu} & \textbf{Mô tả}\\\hline
  \endfirsthead
  \hline\rowcolor{sgCream}
  \textbf{STT} & \normalfont\textbf{Tên trường} & \normalfont\textbf{Kiểu dữ liệu} & \textbf{Mô tả}\\\hline
  \endhead
}{%
  \end{longtable}\normalsize}
